\documentclass[../../main2.tex]{subfiles}

\begin{document}

	% ======================================================================
	% 骨架文件：小节标题与追问链为终版；正文由后续会话填充。
	% 扩增模式标注：本章 = 有压（强制执行）+ 定量编码。
	% 原书材料接入点：
	%   - 原书 part3 ch4（法律）→ §12.1 变长编码 + §12.2 定量意义拆分接入
	%   - §12.2 是全书两半的接合点：ch7 的 C_A 公式 = 法律定责公式
	%   - §12.4 嵌套关系必须用集合符号写出，且必须是推导出来的
	% ======================================================================

	\chapter{通道四：法律（博弈结果的编码与定量方法）}
	\label{ch:law}

	\begin{motivation}
		上一章遗留的问题：权力需要被约束——需要「即使权力方想作恶也无法」的机制。更普遍地：动态博弈产生了结果，双方对结果有争议，谁说了算？\\
		\textbf{核心动机}：论证法律的双重身份——它既是通道（变长编码的规则系统），更是\textbf{第7章贡献分解的现实应用}：$C_A = P(D \mid A) - P(D \mid \neg A)$ 就是法律定责的公式。这一章是全书两半的接合点。\\
		\textbf{扩增模式}：有压——强制执行消耗资源；编码一旦确立，本身抗变化。
	\end{motivation}

	\begin{logicmap}
	\begin{tikzpicture}[node distance=0.9cm, every node/.style={draw, rectangle, rounded corners, align=center}]
	\node (q) {博弈结果有争议，谁说了算？};
	\node (e) [below=of q] {法律是变长编码（§12.1）};
	\node (m) [below=of e] {法律是 Part III 的现实应用（§12.2）};
	\node (d) [below=of m] {法律驱动其余三通道（§12.3）};
	\node (n) [below=of d] {四通道嵌套关系（§12.4）};
	\draw[-{Stealth}] (q) -- (e);
	\draw[-{Stealth}] (e) -- (m);
	\draw[-{Stealth}] (m) -- (d);
	\draw[-{Stealth}] (d) -- (n);
	\end{tikzpicture}
	\end{logicmap}

	% ==================================================================
	\section{法律是变长编码}
	\label{sec:law-encoding}
	为什么法律条文那么长？因为法律是变长编码。

	信息论里，高频事件用短码字，低频事件用长码字——这是霍夫曼编码的基本原理。法律用的是同一个逻辑：常见情形用短规则（「不得杀人」四个字覆盖了绝大多数案件），罕见情形用长规则（几十页的例外条款、但书、限定条件）。歧义的成本高于编码的成本：如果一条规则太短，边界情形就会落入灰色地带，而灰色地带的代价（误判、不公、诉讼）远超多写几页条文的成本。所以民法典必须厚——它不是在堆砌废话，而是在为所有可能的边界情形预留码字。

	「法律不是道德，是道德高频部分的压缩包。」这句话的意思是：法律不试图覆盖所有道德判断（那需要无限长的编码），而是把社会中反复出现、争议最大的那些行为模式提取出来，压缩成可操作的规则。它和语言的关系同构：第9章说的语言压缩层——把高频语义模式固化为词汇和句法——和法律做的是同一件事，只是作用对象不同。

	\begin{readingnote}
		「不得杀人」四个字人人知道，但死刑复核为什么要几百页卷宗？因为那四个字是短码字，覆盖的是「明显故意杀人」这个高频情形。几百页卷宗是在处理长码字部分：他是不是真的「故意」？有没有防卫过当？精神鉴定怎么说？每个边界情形都需要额外的编码来消歧。
	\end{readingnote}

	\paragraph*{逻辑衔接}
	编码讲了法律的「形」。但法律的「魂」不在编码——在它被迫给出的那个数字。下一节是全书两半的接合点。
	% ==================================================================
	\section{法律是 Part III 的现实应用}
	\label{sec:law-quantitative}
	本节是全书两半的接合点。第7章建立的公式 $C_A = P(D \mid A) - P(D \mid \neg A)$——行为 A 对结果 D 的贡献等于「有 A 时 D 的概率」减去「没有 A 时 D 的概率」——就是法律定责的数学内核。

	法律不是「规则体系」，是被强制要求给出确定数字的系统。纠纷必须了结：法官不能说「有点像他的错」，必须说「70% 他的责任，30% 对方的过失」。这个强制给数字的压力，迫使法律发展出社会科学里最精细的因果分解工具。从罗马法的「故意/过失」二分，到现代侵权法的责任比例（comparative negligence），再到产品责任中的市场份额规则——每一步都是在把模糊的道德判断逼向可量化的条件概率差。

	一旦你接受「责任 $\propto P(\text{损害} \mid \text{行为}) - P(\text{损害} \mid \neg \text{行为})$」，你就得到了社会科学因果推断的统一语言。流行病学里的归因风险（attributable risk）、经济学里的反事实分析、政治学里的因果效应估计——都是同一个公式在不同领域的变体。法律是第一个被迫把这个公式操作化的学科，因为它必须给出判决，而判决必须是一个数字。

	\begin{questionbox}
		「法律定责是拍脑袋，哪有真算概率？」显式概率化确实罕见——法官不会在判决书里写「经计算，被告的过错概率为 0.73」。但举证责任和优势证据标准的实质就是门槛化的概率语言：「谁主张谁举证」是在分配先验；「优势证据标准」（preponderance of evidence，即 >50% 可能性）就是在说 $P(D|A) > P(D|\neg A)$。粗粒化，但结构相同。
	\end{questionbox}

	\begin{reader_anticipation}
		\item 条件概率定责的严格形式？→ 第7章 §7.1--7.3 与附录 A.2/A.3。
	\end{reader_anticipation}

	\paragraph*{逻辑衔接}
	法律的定量内核亮出来了。但作为通道，它对其他三条通道做什么？下一节回答。
	% ==================================================================
	\section{法律如何驱动其余三通道}
	\label{sec:law-drives}
	法律改变其他三条通道的方式各有不同：

	对模因通道：法律划定言论的合法性边界。什么话能说、什么话不能说，由法律决定。这不是压制传播，而是给传播加了约束——就像经济通道给模因加了稀缺约束一样。言论自由不是「说什么都可以」，而是「除了法律明确禁止的之外都可以说」。

	对经济通道：法律界定产权。什么是你的、什么是我的、什么可以交易、什么不可以——这些都不是自然事实，而是法律建构。没有产权界定，经济交换就没有基础：你没法卖一件「不确定是不是你的东西」。

	对政治通道：法律划定权力的边界。违宪审查的核心问题就是「权力能不能做这件事」——法律在这里充当了权力的外部约束。第11章说权力会自我扩张，法律就是那个「即使权力方想作恶也无法」的机制。

	这里有一个必须诚实说清的双重身份：法律既是第四通道（它自己产生新的行为——立法、司法、执法），又是前三通道的约束（它限制模因、规范经济、约束政治）。两种身份不矛盾：约束面朝存量（管理已有的行为），通道面朝增量（创造新的规则和制度）。预测时两种用法都有效——分析一个政策的影响时，既要看它作为约束如何改变前三通道的行为，也要看它作为通道如何产生新的博弈。

	\paragraph*{逻辑衔接}
	四条通道全部登场。它们到底是什么关系？下一节给出收口。
	% ==================================================================
	\section{四通道的嵌套关系}
	\label{sec:law-nesting}
	四条通道的关系可以用集合嵌套表达：文化 $\supset$ 经济 $\supset$ 政治。

	文化是最外层：一切社会行为都涉及模因传播，这是基底层。经济是加了稀缺约束的子集：不是所有传播都涉及资源转移，但所有经济行为都涉及传播（信息、信号、预期）。政治是再加了身份权重的子集：不是所有交换都涉及身份关系，但所有政治行为都涉及交换（利益、权力、认同）。法律不是第五层嵌套，而是横跨三层的定量框架——它给每一层提供「怎么算」的工具。

	嵌套意味着中间层不能独立存在：没有不传播模因的经济（交易需要信息），没有不交换的政治（权力需要资源基础），没有不传播的政治（认同需要扩散）。每层「加入的约束」回放：无压 → +稀缺（经济）→ +身份（政治）。

	介观统一论的三个命题在此陈述（只描述，不证明）：经济是稀缺化的模因；政治是人情化的经济；法律是博弈结果的条件概率编码。\lowfreq{三个命题的严格推导见附录/待发表论文《介观统一论》，此处只作直觉呈现。}

	\paragraph*{逻辑衔接}
	分解完成了：四条通道、每条的定义、判据、贡献算法、彼此的嵌套。Part V 回答最后的问题：怎么把贡献表拼回一条时间线。
	% ==================================================================
	\section{本章小结与衔接}
	\label{sec:law-summary}
	Part IV 收束。四条通道——模因、经济、政治、法律——是同一套「认同 + 资源 + 传播」动力学在不同约束下的同一张脸。模因是基底层（无压复制），经济加了稀缺（守恒律），政治加了身份（认同权重），法律提供了定量框架（条件概率编码）。它们不是四个独立的东西，而是一个系统的四个切面。

	\paragraph*{逻辑衔接}
	第13章「从贡献表到未来走向」：拿着 $C_{\text{模因}}, C_{\text{经济}}, C_{\text{政治}}, C_{\text{法律}}$ 四个数字，说出「接下来会怎样」。
\end{document}
